\documentclass[a4paper,12pt]{article}

\usepackage[T2A]{fontenc}
\usepackage[utf8]{inputenc}
\usepackage[russian]{babel}
\usepackage{amsmath,amssymb,mathtools}
\usepackage{PTSans}
\renewcommand{\familydefault}{\sfdefault}
\usepackage{icomma}
\usepackage[
  a4paper,
  left=20mm,
  right=20mm,
  top=20mm,
  bottom=22mm
]{geometry}
\usepackage[nopatch=footnote]{microtype}
\usepackage{tikz}
\usetikzlibrary{calc,arrows.meta,positioning,shapes.geometric}
\usepackage{tabularx,booktabs}
\usepackage{enumitem}
\usepackage{hyperref}
\usepackage{parskip}
\usepackage{fancyhdr}
\usepackage{setspace}
\usepackage{xcolor}

\definecolor{accent}{HTML}{008080}
\definecolor{darkaccent}{HTML}{004D4D}
\definecolor{textgray}{HTML}{333333}
\definecolor{softgray}{HTML}{6B7280}

\hypersetup{
  colorlinks=true,
  linkcolor=darkaccent,
  urlcolor=darkaccent,
  pdftitle={Чек-лист: показательные уравнения и задачи прикладного характера},
  pdfauthor={Лёвин Артём Александрович}
}

\onehalfspacing
\setlength{\parindent}{0pt}
\setlength{\parskip}{0.55em}
\setlist[itemize]{leftmargin=1.6em,itemsep=0.25em,topsep=0.3em}
\setlist[enumerate]{leftmargin=1.9em,itemsep=0.35em,topsep=0.3em}
\emergencystretch=2em
\raggedbottom

\pagestyle{fancy}
\fancyhf{}
\fancyhead[L]{\small\color{softgray}ЕГЭ. Профильная математика}
\fancyhead[R]{\small\color{softgray}Показательные уравнения}
\fancyfoot[C]{\small\color{softgray}\thepage}
\renewcommand{\headrulewidth}{0.35pt}
\renewcommand{\footrulewidth}{0pt}
\setlength{\headheight}{15pt}

\newcommand{\const}{\mathrm{const}}
\newcommand{\keyidea}[1]{\textbf{\color{darkaccent}#1}}
\newcommand{\lessonsection}[1]{%
  \vspace{0.35em}%
  {\Large\bfseries\color{darkaccent}#1\par}%
  \vspace{0.22em}%
  {\color{accent!55}\hrule height 0.55pt}%
  \vspace{0.6em}%
}

\tikzset{
  flow/startstop/.style={
    ellipse,
    draw=darkaccent,
    line width=0.9pt,
    align=center,
    minimum width=36mm,
    minimum height=10mm,
    fill=accent!4,
    font=\small
  },
  flow/action/.style={
    rounded corners=3pt,
    draw=accent!85!black,
    line width=0.8pt,
    align=center,
    text width=100mm,
    minimum height=11mm,
    fill=accent!3,
    font=\small
  },
  flow/decision/.style={
    diamond,
    aspect=2.05,
    draw=darkaccent,
    line width=0.9pt,
    align=center,
    text width=42mm,
    inner sep=1.4pt,
    fill=accent!3,
    font=\small
  },
  flow/side/.style={
    rounded corners=3pt,
    draw=softgray,
    line width=0.8pt,
    align=center,
    text width=43mm,
    minimum height=12mm,
    fill=black!2,
    font=\small
  },
  flow/arrow/.style={
    -{Stealth[length=3mm,width=2mm]},
    line width=0.9pt,
    draw=textgray
  },
  flow/label/.style={
    font=\small\bfseries,
    text=darkaccent,
    fill=white,
    inner sep=1.5pt
  }
}

\begin{document}
\color{textgray}
\hypersetup{pageanchor=false}

% ============================================================
% Титульный лист
% ============================================================

\begin{titlepage}
\thispagestyle{empty}
\centering
\vspace*{22mm}

{\small\bfseries\color{softgray}
ЕГЭ \textbullet{} ПРОФИЛЬНАЯ МАТЕМАТИКА\par}

\vspace{16mm}

{\Huge\bfseries\color{darkaccent}Чек-лист\par}

\vspace{5mm}

{\LARGE\bfseries
Показательные уравнения\par}

\vspace{2mm}

{\Large
и задачи прикладного характера\par}

\vspace{18mm}

{\large
Матвей Горбачев\par}

\vspace{4mm}

{\normalsize
Ключевые методы, преобразования и контроль типичных ошибок\par}

\vfill

{\normalsize
\today\par}

\vspace{7mm}

{\normalsize
Лёвин Артём Александрович\par}

{\normalsize\bfseries
эксклюзивно для Матвея Горбачева\par}

\vspace{24mm}

\begin{tikzpicture}[remember picture,overlay]
  \draw[accent!45,line width=1.05pt]
    ([xshift=20mm,yshift=15mm]current page.south west)
    .. controls +(42mm,15mm) and +(-42mm,11mm) ..
    ([xshift=-20mm,yshift=19mm]current page.south east);

  \draw[darkaccent!22,line width=0.55pt]
    ([xshift=35mm,yshift=11mm]current page.south west)
    .. controls +(34mm,8mm) and +(-34mm,5mm) ..
    ([xshift=-35mm,yshift=12mm]current page.south east);
\end{tikzpicture}

\end{titlepage}

\hypersetup{pageanchor=true}
\pagenumbering{arabic}

% ============================================================
% Основной чек-лист
% ============================================================

\lessonsection{Основной чек-лист}

\keyidea{Цель:}
уверенно решать показательные уравнения тестовой части ЕГЭ и переносить те же преобразования в прикладные задачи со степенными формулами.

\lessonsection{1. Степени: что должно срабатывать автоматически}

Основные правила:
\[
a^m\cdot a^n=a^{m+n},
\qquad
\frac{a^m}{a^n}=a^{m-n}\quad(a\ne0),
\]
\[
(a^m)^n=a^{mn},
\qquad
a^{-n}=\frac1{a^n}\quad(a\ne0).
\]

Для положительного основания
\[
\sqrt[n]{a^m}=a^{m/n}\qquad(a>0).
\]
Если используется корень чётной степени, его подкоренное выражение должно быть неотрицательным.

\textbf{Быстрая самопроверка из домашней работы:}
\[
5^3\cdot5^{-3}=5^{3-3}=5^0=1.
\]

\keyidea{Рабочий принцип:}
составные числа, дроби и корни по возможности переписывайте как степени одного удобного основания.

\begin{center}
\renewcommand{\arraystretch}{1.3}
\begin{tabularx}{\linewidth}{@{}>{\bfseries}l X X@{}}
\toprule
Форма & Удобная запись & Что важно заметить \\
\midrule
$16$ & $2^4$ & ищем простое основание \\
$\dfrac1{16}$ & $2^{-4}$ & знаменатель даёт отрицательную степень \\
$\sqrt[3]{64}$ & $(2^6)^{1/3}=2^2$ & корень переводим в дробную степень \\
$3{,}2\cdot10^6$ & $32\cdot10^5$ & иногда полезно выровнять степень десяти \\
\bottomrule
\end{tabularx}
\end{center}

\lessonsection{2. Показательные уравнения: привести к одному основанию}

Если
\[
a^{F(x)}=a^{G(x)},
\qquad a>0,
\qquad a\ne1,
\]
то
\[
F(x)=G(x).
\]

\textbf{Самый простой пример:}
\[
2^x=16=2^4
\quad\Longrightarrow\quad
\boxed{x=4}.
\]

\textbf{Пример с дробью в основании:}
\[
\left(\frac12\right)^{6-2x}=4.
\]
Переходим к основанию $2$:
\[
2^{-6+2x}=2^2.
\]
Следовательно,
\[
-6+2x=2,
\qquad
\boxed{x=4}.
\]

\textbf{Пример с корнем и знаменателем:}
\[
8^x=\frac1{\sqrt[3]{64}}.
\]
Так как
\[
8=2^3,
\qquad
\frac1{\sqrt[3]{64}}
=\frac1{2^2}=2^{-2},
\]
получаем
\[
2^{3x}=2^{-2},
\qquad
3x=-2,
\qquad
\boxed{x=-\frac23}.
\]

\keyidea{Контроль знака:}
если степень оказалась в знаменателе, показатель меняет знак.

\textbf{Пример с дополнительным множителем:}
\[
3\cdot9^{x-1}=\frac1{27}.
\]
Все части выражаем через основание $3$:
\[
3^1\cdot(3^2)^{x-1}=3^{-3},
\]
\[
3^{1+2x-2}=3^{-3}.
\]
Поэтому
\[
2x-1=-3,
\qquad
\boxed{x=-1}.
\]

% ============================================================
% Блок-схема 1
% ============================================================

\clearpage
\thispagestyle{empty}

\begin{center}
{\Large\bfseries\color{darkaccent}
Алгоритм: показательное уравнение через общее основание}
\end{center}

\vspace{2mm}

\begin{center}
\begin{tikzpicture}[node distance=7mm and 14mm]

\node[flow/startstop] (start) {Начало};

\node[flow/action,below=of start] (inspect)
{Проверить запись:\\дроби, корни, отрицательные показатели, множители};

\node[flow/action,below=of inspect] (convert)
{Представить числа и корни как степени\\удобного простого основания};

\node[flow/decision,below=9mm of convert] (samebase)
{Получилось\\одно основание\\в обеих частях?};

\node[flow/startstop,right=18mm of samebase,text width=40mm] (other)
{Перейти\\к другому методу};

\node[flow/action,below=9mm of samebase] (collect)
{Собрать каждую сторону в одну степень: $a^{F(x)}=a^{G(x)}$};

\node[flow/action,below=of collect] (equal)
{Приравнять показатели: $F(x)=G(x)$};

\node[flow/action,below=of equal] (solve)
{Решить полученное уравнение};

\node[flow/action,below=of solve] (check)
{Проверить знаки, ограничения исходного выражения\\и найденное значение};

\node[flow/startstop,below=of check] (finish) {Ответ};

\draw[flow/arrow] (start) -- (inspect);
\draw[flow/arrow] (inspect) -- (convert);
\draw[flow/arrow] (convert) -- (samebase);
\draw[flow/arrow] (samebase) -- node[flow/label,left]{да} (collect);
\draw[flow/arrow] (samebase) -- node[flow/label,above]{нет} (other);
\draw[flow/arrow] (collect) -- (equal);
\draw[flow/arrow] (equal) -- (solve);
\draw[flow/arrow] (solve) -- (check);
\draw[flow/arrow] (check) -- (finish);

\end{tikzpicture}
\end{center}

\clearpage

\textbf{Пример для самостоятельного контроля преобразований:}
\[
\frac1{\sqrt[5]{2}}\cdot2^{3x-1}=\sqrt[4]{32}.
\]
Переходим к степеням двойки:
\[
2^{-1/5}\cdot2^{3x-1}=2^{5/4},
\]
\[
2^{3x-6/5}=2^{5/4}.
\]
Следовательно,
\[
3x-\frac65=\frac54,
\qquad
3x=\frac54+\frac65
=\frac{25}{20}+\frac{24}{20}
=\frac{49}{20},
\qquad
\boxed{x=\frac{49}{60}}.
\]

% ============================================================
% Дробные степени неизвестной величины
% ============================================================

\lessonsection{3. Когда сама неизвестная стоит в дробной или целой степени}

В уравнении
\[
x^r=A
\]
цель состоит в том, чтобы получить первую степень неизвестной. Однако преобразование должно быть равносильным на ОДЗ исходного выражения.

\textbf{Пример 1.}
\[
x^{3/5}=\sqrt[5]{2}.
\]
Правая часть положительна, поэтому искомое $x$ также положительно. Можно возвести обе части в степень $5/3$:
\[
x=\left(2^{1/5}\right)^{5/3}
=2^{1/3}.
\]
Итак,
\[
\boxed{x=\sqrt[3]{2}}.
\]

\textbf{Пример 2. Чётная степень требует контроля знака.}
\[
t^4=2^8\cdot10^{12}.
\]
Поскольку
\[
2^8\cdot10^{12}
=\left(2^2\cdot10^3\right)^4
=4000^4,
\]
чисто алгебраическое уравнение имеет два решения:
\[
\boxed{t=\pm4000}.
\]
Если по смыслу задачи $t$ обозначает положительную величину, остаётся значение $t=4000$.

\textbf{Пример 3. Положительная прикладная величина.}
Пусть $V>0$ и
\[
V^{2/3}=16.
\]
Возводим обе части в степень $3/2$:
\[
V=16^{3/2}=(4^2)^{3/2}=4^3,
\]
\[
\boxed{V=64}.
\]

\keyidea{Перед использованием обратной степени проверьте:}
\begin{itemize}
  \item допускается ли такое преобразование на ОДЗ;
  \item может ли чётная степень дать два знака;
  \item задана ли положительность величины смыслом задачи.
\end{itemize}

% ============================================================
% Прикладные задачи
% ============================================================

\lessonsection{4. Прикладные задачи: перевод условия в вычисления}

В задачах этого типа формула уже дана. Основная работа --- аккуратно прочитать обозначения, подставить данные и свести вычисления к знакомой работе со степенями.

\begin{enumerate}
  \item Выпишите формулу, известные величины и то, что требуется найти.
  \item Проверьте согласованность единиц измерения.
  \item Подставляемое выражение берите в скобки: если $V_1=2V_2$, то $V_1^a=(2V_2)^a$.
  \item Десятичный показатель удобно переводить в обыкновенную дробь: $1{,}4=\dfrac75$.
  \item После получения ответа проверьте его знак и физический смысл.
\end{enumerate}

Запись $m(t)$ означает значение функции $m$ при аргументе $t$. Это не произведение $m\cdot t$.

\lessonsection{5. Тип 1: прямая подстановка в степенную формулу}

Пусть
\[
pV^{5/3}=10^5,
\qquad
p=3{,}2\cdot10^6\ \text{Па}.
\]
Найдём $V$.

Удобно переписать
\[
3{,}2\cdot10^6=32\cdot10^5.
\]
Тогда
\[
32\cdot10^5\,V^{5/3}=10^5,
\]
\[
V^{5/3}=\frac1{32}=2^{-5}.
\]
Объём положителен, поэтому
\[
V=\left(2^{-5}\right)^{3/5}=2^{-3}=\frac18.
\]
Итак,
\[
\boxed{V=0{,}125\ \text{м}^3}.
\]

\lessonsection{6. Тип 2: закон радиоактивного распада}

Модель имеет вид
\[
m(t)=m_0\cdot2^{-t/T},
\]
где $m_0$ --- начальная масса, $T$ --- период полураспада.

Пусть
\[
m_0=40\ \text{мг},
\qquad
T=10\ \text{мин},
\qquad
m(t)=5\ \text{мг}.
\]
Тогда
\[
5=40\cdot2^{-t/10},
\]
\[
\frac18=2^{-t/10}.
\]
Так как
\[
\frac18=2^{-3},
\]
получаем
\[
-\frac{t}{10}=-3,
\qquad
\boxed{t=30\ \text{мин}}.
\]

\lessonsection{7. Тип 3: два состояния и формула $pV^a=\const$}

Если в процессе величина
\[
pV^a
\]
остаётся постоянной, то для двух состояний системы
\[
p_1V_1^a=p_2V_2^a.
\]

\textbf{Пример 1. Найти наименьший показатель $a$.}

Объём уменьшили в $2$ раза:
\[
V_1=2V_2.
\]
Давление должно увеличиться не менее чем в $4$ раза:
\[
p_2\ge4p_1.
\]
Из закона процесса
\[
\frac{p_2}{p_1}
=\left(\frac{V_1}{V_2}\right)^a
=2^a.
\]
Следовательно,
\[
2^a\ge4=2^2.
\]
Так как функция $2^a$ возрастает,
\[
a\ge2.
\]
Поэтому
\[
\boxed{a_{\min}=2}.
\]

\keyidea{Почему можно брать границу:}
при поиске наименьшего $a$ условие ``не менее'' достигается впервые при равенстве.

\textbf{Пример 2. Найти новый объём.}

Пусть
\[
pV^{1{,}4}=\const,
\qquad
p_1=1\ \text{атм},
\qquad
V_1=1{,}6\ \text{л},
\qquad
p_2=128\ \text{атм}.
\]
Тогда
\[
p_1V_1^{7/5}=p_2V_2^{7/5},
\]
\[
1{,}6^{7/5}=128\,V_2^{7/5}.
\]
Отсюда
\[
V_2^{7/5}=\frac{1{,}6^{7/5}}{128}.
\]
Поскольку $V_2>0$, возводим обе части в степень $5/7$:
\[
V_2
=\frac{1{,}6}{128^{5/7}}.
\]
Так как
\[
128=2^7,
\qquad
128^{5/7}=2^5=32,
\]
получаем
\[
\boxed{V_2=0{,}05\ \text{л}}.
\]

% ============================================================
% Блок-схема 2
% ============================================================

\clearpage
\thispagestyle{empty}

\begin{center}
{\Large\bfseries\color{darkaccent}
Алгоритм: прикладная задача со степенной формулой}
\end{center}

\vspace{4mm}

\begin{center}
\begin{tikzpicture}[node distance=6mm and 13mm]

\node[flow/startstop] (start2) {Начало};

\node[flow/action,below=of start2] (read)
{Выписать формулу, данные, неизвестную\\и единицы измерения};

\node[flow/decision,below=7mm of read] (twostate)
{Есть два состояния\\одной системы?};

\node[flow/side,right=15mm of twostate] (direct)
{Подставить данные\\в исходную формулу};

\node[flow/action,below=7mm of twostate] (states)
{Записать равенство состояний,\\например $p_1V_1^a=p_2V_2^a$};

\node[flow/action,below=of states] (substitute)
{Перевести условие в формулы и подставить:\\$V_1=kV_2$, $p_2\ge rp_1$ и т.\,п.};

\node[flow/action,below=of substitute] (power)
{Упростить степени и решить полученное\\степенное или показательное уравнение};

\node[flow/decision,below=7mm of power] (extreme)
{Нужно найти\\минимум или максимум?};

\node[flow/side,right=15mm of extreme] (boundary)
{Перейти к граничному\\равенству условия};

\node[flow/action,below=7mm of extreme] (units)
{Проверить ОДЗ, знак, смысл результата\\и единицы измерения};

\node[flow/startstop,below=of units] (finish2) {Ответ};

\draw[flow/arrow] (start2) -- (read);
\draw[flow/arrow] (read) -- (twostate);
\draw[flow/arrow] (twostate) -- node[flow/label,left]{да} (states);
\draw[flow/arrow] (twostate) -- node[flow/label,above]{нет} (direct);
\draw[flow/arrow] (direct.south) |- (substitute.east);
\draw[flow/arrow] (states) -- (substitute);
\draw[flow/arrow] (substitute) -- (power);
\draw[flow/arrow] (power) -- (extreme);
\draw[flow/arrow] (extreme) -- node[flow/label,left]{нет} (units);
\draw[flow/arrow] (extreme) -- node[flow/label,above]{да} (boundary);
\draw[flow/arrow] (boundary.south) |- (units.east);
\draw[flow/arrow] (units) -- (finish2);

\end{tikzpicture}
\end{center}

\clearpage

% ============================================================
% Ошибки и самопроверка
% ============================================================

\lessonsection{8. Типичные ошибки и быстрая самопроверка}

\begin{enumerate}
  \item \textbf{Знак в знаменателе.}
  \[
  \frac1{2^5}=2^{-5}.
  \]

  \item \textbf{Степени одного основания.}
  При умножении показатели складываются; при делении --- вычитаются.

  \item \textbf{Подстановка без скобок.}
  Если $V_1=2V_2$, то
  \[
  V_1^a=(2V_2)^a,
  \]
  а не $2V_2^a$.

  \item \textbf{Десятичная степень.}
  Для вычислений
  \[
  1{,}4=\frac75,
  \qquad
  \frac1{1{,}4}=\frac57.
  \]

  \item \textbf{Чётная степень.}
  Из равенства $t^4=A>0$ в алгебраической задаче обычно следуют два значения $t=\pm\sqrt[4]{A}$; положительный знак оставляют только при соответствующем ограничении.

  \item \textbf{Слова ``не менее'' и ``не более''.}
  Они задают неравенство. Равенство используют на границе только тогда, когда задача просит минимальное или максимальное допустимое значение.

  \item \textbf{Сокращение.}
  Перед сокращением выражения, зависящего от переменной, нужно убедиться, что оно не равно нулю. В формулах газа $p>0$ и $V>0$ следуют из физического смысла.
\end{enumerate}

\keyidea{Контроль перед ответом:}
проверьте основание степени, знак показателя, ОДЗ, подстановку в скобках и единицы измерения.

% ============================================================
% Тренировочный блок
% ============================================================

\clearpage

\lessonsection{Тренировочный блок}

\textbf{Путь А.} Выполните $5$ заданий. Решения оформляйте последовательно; в прикладных задачах сначала выписывайте формулу и данные.

\begin{enumerate}[label=\textbf{\arabic*.}]

  \item Решите уравнение
  \[
  \left(\frac12\right)^{8-2x}=4.
  \]

  \item Решите уравнение
  \[
  \frac1{\sqrt[3]{2}}\cdot2^{2x+1}=\sqrt[6]{4}.
  \]

  \item Величины $p$ и $V$ связаны формулой
  \[
  pV^{3/2}=64.
  \]
  Найдите положительный объём $V$, если $p=8$ в согласованных единицах.

  \item Масса вещества меняется по закону
  \[
  m(t)=m_0\cdot2^{-t/T}.
  \]
  В начальный момент $m_0=96\ \text{мг}$, период полураспада $T=12\ \text{мин}$. Через сколько минут масса станет равной $12\ \text{мг}$?

  \item Для газа выполняется закон
  \[
  pV^a=\const,
  \qquad a>0.
  \]
  Объём газа уменьшили в $16$ раз. Найдите наименьшее значение $a$, при котором давление увеличится не менее чем в $8$ раз.

\end{enumerate}

% ============================================================
% Ответы
% ============================================================

\clearpage

\lessonsection{Ответы}

\begin{center}
\renewcommand{\arraystretch}{1.45}
\begin{tabularx}{0.82\linewidth}{@{}c X@{}}
\toprule
\textbf{№} & \textbf{Ответ} \\
\midrule
1 & $x=5$ \\
2 & $x=-\dfrac16$ \\
3 & $V=4$ \\
4 & $36\ \text{мин}$ \\
5 & $a_{\min}=\dfrac34$ \\
\bottomrule
\end{tabularx}
\end{center}

\vfill

{\small\color{softgray}
Главный ориентир: приводите выражения к удобной степенной форме, а затем отдельно контролируйте ограничения, знаки и смысл переменных.
}

\end{document}